%======================================================================
%  Extended Data Figures
%
%  Five figures promoted out of the Supplementary set at the author's request
%  (Nature Biotechnology displays Extended Data alongside the main figures
%  online). Order follows their order of appearance in the manuscript:
%
%    ED 1  <- former Supp. Fig. S1   IMG/PR collection, hierarchy and sampling
%    ED 2  <- former Supp. Fig. S3   window -> IGR origin-inference pipeline
%    ED 3  <- former Supp. Fig. S6   additional window-probability tracks
%    ED 4  <- former Supp. Fig. S7   probe layer selection
%    ED 5  <- former Supp. Fig. S17  embedding projection by annotation
%
%  Captions are unchanged from their Supplementary versions -- the GPN-Star
%  convention is identical for both bands (no bold lead sentence, opening
%  directly with the description). Only the \label prefixes changed
%  (fig:supp-* -> fig:ed-*) so the cross-reference wrappers read correctly.
%
%  This section sits BEFORE \maintextpartfalse, so anything cited here counts
%  as a main-text citation. None of these five captions carries a \cite, so
%  the two reference lists are unaffected.
%
%  Snapshot 2026-08-28; promoted 2026-09-01.
%======================================================================

% ---- ED 1: corpus, hierarchy and diversity-aware sampling (was Supp. S1)
\begin{figure}[h]
    \centering
    \fitfig{figures/figS1_combined}
\caption{
The IMG/PR collection, corpus construction, cluster hierarchy, and diversity-aware sampling. Panels \textbf{a}, \textbf{c} and \textbf{d} describe the full IMG/PR collection (699,973 sequences, 214,950 PTUs); the remaining panels describe the 693,638-sequence training corpus after quality filtering and the two-species holdout. Panel titles are given here rather than on the artwork. \textbf{a}, Host phyla by plasmid-sequence count (log scale) --- the twelve largest plus an aggregate "other phyla" bar (49 phyla in all; 40\% of sequences carry a host assignment). \textbf{b}, Sequences removed at each filtering step (log scale); 99.1\% of the collection is retained. \textbf{c}, Concentration of diversity: cumulative fraction of dereplicated sequences (solid) and of sub-PTUs (dashed) against the cumulative fraction of PTUs, ranked and shaded by category (Major / Moderate / Unique); a small minority of PTUs accounts for most of the collection. \textbf{d}, Length distribution of the full collection (log-scaled); the dashed lines mark the 2,048 bp minimum and the 10.1 kb median, and the sub-2,048 bp sequences the filter removes are visible to the left. \textbf{e}, Completeness (complete vs incomplete) and topology (linear / direct- or inverted-terminal-repeat) annotations. \textbf{f}, Number of PTUs against sub-PTUs per PTU (both log); 81\% of PTUs resolve to a single sub-PTU, the largest to 735. \textbf{g}, Cluster-size distributions at the two identity levels --- dereplication (99\% ANI) and sub-PTU (95\% ANI). \textbf{h}, Distribution of expected Major-PTU allocations (mean 1.42, range 1.08--3.15, none at the floor). \textbf{i}, The Major-PTU allocation $a_i$ against sub-PTU count $n_i$: it grows sublinearly as $n^{2/3}$ (dashed), floored at one representative and capped at the 95th percentile of sub-PTU count ($Q_{0.95} = 10$ sub-PTUs), giving an expected-allocation ceiling of 3.15 representatives per epoch.
}
\label{fig:ed-corpus}
\end{figure}

% ---- ED 2: window -> intergenic-region inference pipeline (was Supp. S3)
%      Raster-only in the source artifact; artwork copied in as-is.
%
%      2026-09-01, author-directed: caption retitled. It opened "Genome-wide
%      origin discovery with GPN-Plasmids", which named the downstream
%      application rather than the artwork. The figure shows ONE query
%      plasmid going from window-level classification to intergenic-region
%      classification, and the two things that buys -- localization within a
%      region, and a ranking over a plasmid's regions. Genome-wide
%      application is a consequence of the pipeline, so it is demoted to the
%      closing clause. The figure is now cited from §3.1, which describes
%      the same two stages in prose.
\begin{figure}[h]
    \centering
    \fitfig{figures/supp_inference_pipeline.png}
\caption{
From window-level classification to intergenic-region ranking with \model. A query plasmid is annotated into coding sequences and intergenic regions (IGRs), and each IGR is tiled into overlapping 512 bp windows that frozen \model embeddings and the trained probe score for origin overlap. Reading those per-window scores along the region gives a calibrated probability profile that localizes origin signal to part of an IGR rather than only flagging the IGR as a whole. Pooling the profile at the 75th percentile (or taking the maximum if the IGR has five or fewer windows) and converting it per split into one ensemble score per IGR reduces the region to a single value, so a plasmid's IGRs become comparable to one another; ranking them by that score gives its candidate origin-containing regions. The same two-stage workflow, applied unchanged across IMG/PR, is the source of the genome-wide predictions used downstream (\secref{sec:genomewide}).
}
\label{fig:ed-pipeline}
\end{figure}

% ---- ED 3: additional window-probability tracks (was Supp. S6)
\begin{figure}[h]
    \centering
    \fitfig{figures/supp_region_tracks_8}
\caption{
Calibrated window-probability tracks for additional origin-containing regions. Eight held-out examples --- four \emph{oriT} (\textbf{a}--\textbf{d}, left) and four \emph{oriV} (\textbf{e}--\textbf{h}, right) --- scored by the \model probe: calibrated per-window probability (red, left axis) along a 512 bp window stepped at 16 bp, with the annotated origin shaded orange, flanking coding sequence grey, and the per-window origin overlap (blue dashed, right axis). The green line is the 75th-percentile region score used for ranking. Same tracks and legend as \fref{fig:origin}e.
}
\label{fig:ed-tracks8}
\end{figure}

% ---- ED 4: probe layer selection by validation AUPRC (was Supp. S7)
\begin{figure}[h]
    \centering
    \fitfig{figures/supp_win_layer}
\caption{
Probe layer selection by validation AUPRC. Window-level AUPRC for every probed layer of each pretrained embedder (per-split mean bars over per-split heatmaps, best-val layer outlined); \textbf{a}, \emph{oriT}; \textbf{b}, \emph{oriV}. Each embedder's reported probe uses the single layer maximizing mean validation AUPRC across both tasks --- \model \texttt{blocks.60}, Evo2-1B/7B \texttt{blocks.15.mlp.l3}, Evo2-40B \texttt{blocks.16.mlp.l3}, PlasmidGPT \texttt{h.3}. Selection uses validation columns only; because parts rotate, each part is validation in one split and test in another. GPN reporting considered layers 60--72.
}
\label{fig:ed-layer}
\end{figure}

% ---- ED 5: embedding projection by coarse annotation (was Supp. S17)
\begin{figure}[h]
    \centering
    \fitfig{figures/figS5_umap_feature}
\caption{
Embedding projection coloured by coarse functional annotation. Two-dimensional UMAP of frozen \model embeddings over 308,478 non-overlapping 10 bp tiles from 40 origin-enriched Enterobacterales plasmids, each tile embedded from its own 2,048 bp context and averaged over forward and reverse-complement inputs (\nameref{sec:methods}). Each point is one tile, coloured by its annotation: intergenic (blue), coding sequence (orange), origin of vegetative replication (red), origin of transfer (green) and non-coding RNA (purple). Labels are assigned by overlap across the whole tile, and origin calls are curated BLAST hits that replace PlasAnn's own. Coding and intergenic tiles occupy largely distinct regions, and the two origin classes and non-coding RNA are associated with compact regions within the intergenic territory. No separation metric or classifier was computed, so this is an association rather than a demonstrated separation, and the cohort is origin-enriched by construction. Annotations were not used during pretraining or in fitting the projection; they were used to filter ambiguous tiles and to colour the points.
}
\label{fig:ed-umap-feature}
\end{figure}


%======================================================================
%  Extended Data Table 1 -- origin task / training / inference parameters
%
%  Moved here from the main text 2026-09-04. It was main-text Table 3, a
%  32-row longtable of Methods parameters. Nature Biotechnology allows six
%  main display items and the manuscript had seven (Figures 1-4, Tables 1-3);
%  a parameter table is the one that least needs a main-text slot. Extended
%  Data allows ten items and carried only five, so this costs nothing.
%  \beginextendeddata has already reset the table counter, so this typesets as
%  "Extended Data Table 1". The caption lost its pointer to main-text Table 3
%  and its "Design-critical values appear inline" clause, which described the
%  old split between this table and the Results prose; that prose is now in
%  Methods. The \label is unchanged, so \edtref{tab:origin-params} resolves.
%======================================================================

{\small
\begin{longtable}{p{0.40\textwidth}p{0.30\textwidth}p{0.20\textwidth}}
\caption{
Key task, training and inference parameters for the origin classifier. The provenance column distinguishes values set at run time from values fixed in the analysis source and from library defaults. These parameters accompany the corresponding \nameref{sec:methods} subsections.
}\label{tab:origin-params}\\
\toprule
Parameter & Value & Provenance \\
\midrule
\endfirsthead
\multicolumn{3}{l}{\emph{\tablename~\thetable{} (continued)}}\\
\toprule
Parameter & Value & Provenance \\
\midrule
\endhead
\bottomrule
\endfoot
\multicolumn{3}{l}{\textit{Reference curation \& detection}} \\
Curated origins (\emph{oriT} / \emph{oriV}) & 112 / 224 & derived (MMseqs2 dedup) \\
MMseqs2 dedup (id / cov / mode) & 90\% / 90\% / shorter & set at run time \\
BLASTN detection (word size / dust) & 7 / off & set at run time \\
Detection thresholds (identity / subj.\ coverage) & 80\% / 80\% & fixed in source \\
\midrule
\multicolumn{3}{l}{\textit{Regions \& selection}} \\
CDS caller & Pyrodigal-gv (meta) & library default \\
Min intergenic length / min regions per plasmid & 50 bp / 3 & fixed in source \\
Origin--region assignment & $\geq$ 50\% origin length, summed across intersecting regions & fixed in source \\
Plasmids per unique origin / max origins per plasmid & top 5 / $\leq$ 5 & fixed in source \\
\midrule
\multicolumn{3}{l}{\textit{Splits}} \\
Similarity (EMBOSS Water gap open / extend) & 10 / 0.5 & set at run time \\
GraphPart folds / threshold & 5 / 0.3 & set at run time \\
\midrule
\multicolumn{3}{l}{\textit{Windowing \& imbalance}} \\
Window size / training step / inference step & 512 / 256 / 16 bp & set at run time \\
Positive-window rule & $\geq$ 50\% origin or $\geq$ 256 bp & fixed in source \\
Downsample neg:pos (\emph{oriT} / \emph{oriV}) & 20:1 / 20:1 & set at run time \\
Positive-class loss weight (\emph{oriT} / \emph{oriV}) & 20$\times$ / 20$\times$ & set at run time \\
Reverse-complement augmentation & enabled & fixed in source \\
\midrule
\multicolumn{3}{l}{\textit{Probes \& training}} \\
MLP probe input & mean-pooled window & fixed in source \\
One-hot probe (hidden / kernel / dropout) & 64 / 16 / 0.15 & set at run time \\
Stochastic runs $\times$ splits per (embedder, layer) & $5 \times 5 = 25$ & set at run time \\
Optimizer / learning rate / weight decay & AdamW / $5\times10^{-4}$ / 0.01 & set at run time \\
Schedule / epochs / batch size & cosine / 30 / 256 & set at run time \\
Model selection & best val AUPRC & fixed in source \\
Layer selection & mean val AUPRC, both tasks & derived \\
\midrule
\multicolumn{3}{l}{\textit{Calibration \& aggregation}} \\
Calibrator & isotonic (clip) & fixed in source \\
Ensembling & probability soft-vote (5 runs) & fixed in source \\
Region aggregator & percentile-75 & set at run time \\
\end{longtable}
}
